\documentclass[11pt, a4paper]{article}
\usepackage[utf8]{inputenc}
\usepackage[spanish]{babel}
\usepackage{amsmath, amssymb}
\usepackage{geometry}
\usepackage{hyperref}
\usepackage{enumitem}

\geometry{top=2.5cm, bottom=2.5cm, left=2.5cm, right=2.5cm}

\hypersetup{
	colorlinks=true,
	linkcolor=blue,
	urlcolor=blue
}

\title{\Large \textbf{Pasos de instalación}}
\author{}
\date{}

\begin{document}
	
	\maketitle
	\vspace{-1cm}
	
	\noindent Realice los siguientes pasos para .
	
	\begin{enumerate}[leftmargin=*]
 \item Entre al sitio \url{https://julialang.org/downloads/} y descargue la versión más reciente de Julia  para su equipo. Siga las instrucciones dadas en dicha página, para su sistema operativo.Instalelo con un PATH simple y sin espacios por ejemplo ``C:$\backslash$julia''.
 
 Si el programa de instalación sugiere agregar la trayectoria a PATH, acéptelo; en caso contrario, copie el PATH donde se instala Julia para usarla en la configuración posterior.
			

		
\item \textbf{Ahora instale el editor integrado Visual Studio Code VSCode con el cual vamos a trabajar.}
		\begin{enumerate}[label=\alph*)]

\item Descargue e instale el editor (también open source) de la página \url{https://code.visualstudio.com/Download}. Siga las instrucciones para su instalación.

\item Ahora configure VSCode para usar con Julia. Para ello, debe buscar en la opción \textit{Extensions} (columna izquierda de la pantalla, u oprimiendo \texttt{Ctrl+Shift+X}). Se abre una columna de búsqueda llamada \textbf{EXTENSIONS: MARKETPLACE}: teclee \texttt{julia}.\\
			Aparecen varias extensiones, de las cuales debe seleccionar una que dice \textbf{Julia - Julia Language Support}. Cliquee sobre el botón \textit{Install}. Cuando termine de instalar la extensión, reinicie VSCode. Ya debería estar listo para trabajar en Julia.
		\end{enumerate}

\item Instale la extensión para \texttt{Jupyter} en VSCode 


\item Cree una carpeta para control donde va a trabajar durante el semestre en el laboratorio. Trate de que sea una carpeta con un PATH simple, por ejemplo, ``C:$\backslash$labcontrol''. 

\item En el menú principal del VSCode, seleccione la opción \textbf{File $=>$ Open Folder} y abra la carpeta que habia creado en el paso anterior. Cree un archivo con extensión jl, por ejemplo, ``instalar.jl''. Copie las siguientes instrucciones en el archivo:

\verb*|print("Instalando IJulia...")|\\
Using Pkg\\
\verb*|Pkg.add("IJulia")|



\item Instale el \emph{firmware} de la tarjeta, según las instrucciones \href{https://nebisman.github.io/DCMotor.jl/dev/instalacion/firmware/#Instalaci%C3%B3n-del-Firmware}{según las instrucciones de este enlace.}  

\item En el mismo archivo ``instalar.jl'' que creó anteriormente \href{https://nebisman.github.io/DCMotor.jl/dev/instalacion/paquete/}{ corra las instrucciones de instalación del paquete en este enlace.}  

\item Descarque el archivo de la práctica 4 llamado ``practica4.ipynb'' y cópielo en su carpeta de trabajo. En seleccionar kernel (parte superior derecho), seleccione ``jupyter kernel'' y luego julia. Si no aparece,  recargue la ventana con Ctrl+Shift+P → ``Developer: Reload Window'' e intente de nuevo.

\item Verifique que la primera celda corre.


		
%\item Por último, vamos a instalar algunos paquetes que usaremos durante el semestre. Algunos son bastante grandes y pueden tomar un tiempo considerable para instalarlos.
%		
%\item Copie el instalador adjunto (\texttt{instalador\_dos.jl}). Si todo sale bien, ya está.\\
%		En caso de que tengan algún tropiezo para instalar alguna componente, se recomienda mirar el blog: \url{https://blog.glcs.io/install-julia-and-vscode} para la instalación y \url{https://blog.glcs.io/install-julia-and-vscode#heading-configuring-vs-code-for-julia} para la configuración del VSCode.
%		
%		\item Instale la extensión para \texttt{Jupyter} en VSCode.
\end{enumerate}
%	
%\vspace{0.5cm}
%	\noindent Si tienen preguntas adicionales, trataremos de responderlas en la clase.
	
\end{document}